\section{Analyse de stabilité}
\paragraph{Interêt d'une étude de stabilité} Etude de l'effet d'une perturbation sur un état de base. Connaitres les états et les bifurcations possibles. Le souci c'est la complexité du système dynamique étudié. \\

\subsection{Obtention des EDO d'un système dynamique}
Étude mené sur des système dynamiques = EDO
 \begin{equation}
 	\label{eq:EDOdyn}
 	\dot X = F(X, r, t) 
 \end{equation}
 mais les équations physiques, notamment de la dynamique, sous formes d'EDP 
 \begin{equation}
 	\label{eq:EDPphys}
 	\frac{\partial V}{\partial t} = \mathcal L(V) + \mathcal{N} (V)
\end{equation}
 où $\mathcal L$ est un opérateur linéaire. On fait des études autour d'une situation de base connue. $\mathcal{L}$ fait apparaitre les cohérence spatiales, retroactions linéaires dans ses modes propres. \\
 On cherche des solutions à variables séparées et si $\mathcal{L}$ peut etre auto-adjoint c'est chouette :) \\
 Le problème apres c'est plus de traiter $\mathcal{N}$ mais avec dvp de Taylor on peut se ramener à une équatio nde la dynamique. \\
 \textbf{Objectif ? } Est ce que ça peut etre un objectif de mettre le problème sous forme de problème dynamique? \\

\subsection{Cas de la convection de Rayleigh-Bénard}
On s'intéresse au cas de Rayleigh-Bénard avec aucune dépendance en $y$, donc que $\partial_y = 0$ et $v=0$. 
fluide incompressible (expliquer pk?)\\
Système étudié décrit par les équations de Boussinesq : 
\begin{equation}
 	\begin{cases}
 		\underline \nabla \cdot \underline u = 0 \\
 		\left(\partial_t  + (\underline u \cdot \underline \nabla) \right) \underline u = \nu \underline \nabla ^2 \underline u - \frac{1}{\rho_0} \nabla p + \rho g \\
 		\left(\partial_t  + (\underline u \cdot \underline \nabla) \right) \theta = \kappa \underline \nabla ^2 \theta + N^2 w   \\
 	\end{cases}
 \end{equation}
Ce pb est étudié pour Rayleigh-Bénard = instabilité de convection. On trouve un point de bifurcation $\propto \Delta t$, cas pour les rouleaux plus compliqué car instabilité de vent/ cisaillement ? étude par \cite{liu_analytical_2009}.
% Study of stability of the Rayligh Benard convection => Rayleigh number to decide when the convection begins. Rolls is basicely the same but with wind ? so mabe try to do the same 
% Study of unstabilities that lead to rolls= expected basic state of hurricane BL (\cite{foster_why_2005}. )
